\section{Weil classes in coordinates}\label{sec:explicit}

\Cref{sec:weiltype} defines the Weil line as an intersection inside
$H^{2n}(A,\CC)$. That description is the right one for proofs, and it is the
wrong one for getting a feel for the object, because it never exhibits a single
class. This section writes the Weil classes out. We build an abelian variety of
$(K,-1,n)$-Weil type from a $2n$-dimensional $K$-vector space and a hermitian
form, produce two explicit integral generators $\omega_{1},\omega_{2}$ of the
Weil line as sums of $2^{2n-1}$ monomials each, display them in full for $n=2$
and tabulate them for $n=3$, and then use the coordinates for two things the
abstract description does not give: an elementary second proof that the Weil
line misses the divisor subring, and a complete count of the Hodge classes in
every degree.

Everything in this section is verified symbolically in
\texttt{explicit\_weil.py} and \texttt{hodge\_invariants.py}; the checks are
listed in \Cref{sec:verification}, items (VI) and (VII).

\subsection{The standard model}

\begin{setupenv}[The model]\label{setup:model}
Fix a squarefree positive integer $d$, let $K=\QQ(\sqrt{-d})$, and fix the
embedding $K\hookrightarrow\CC$ carrying $\sqrt{-d}$ to $\delta:=i\sqrt{d}$.
Fix $n\ge1$ and let
\[
  V=K^{2n},
  \qquad
  u_{1},\dots,u_{2n}\ \text{the standard $K$-basis},
\]
regarded as a $\QQ$-vector space of dimension $4n$ with the basis
\[
  x_{j}:=u_{j},
  \qquad
  y_{j}:=\sqrt{-d}\,u_{j},
  \qquad 1\le j\le 2n .
\]
Let $S\colon V\to V$ be multiplication by $\sqrt{-d}$, so that in this basis
\begin{equation}\label{eq:Sbasis}
  S x_{j}=y_{j},
  \qquad
  S y_{j}=-d\,x_{j}.
\end{equation}
Let $\epsilon_{j}=-1$ for $1\le j\le n$ and $\epsilon_{j}=+1$ for
$n<j\le 2n$,%
\footnote{The signs are not free. The proof of \Cref{thm:model} shows that the second Riemann relation holds for this choice and fails as soon as one of them is reversed, so the signature $(n,n)$ of the sign vector is the balanced condition in disguise. Readers who prefer the opposite global sign should replace $Q$ by $-Q$ throughout and interchange $V^{1,0}$ with $V^{0,1}$.} and let $Q$ be the alternating $\QQ$-bilinear form on $V$ with
\begin{equation}\label{eq:Qbasis}
  Q(x_{j},y_{j})=\epsilon_{j},
  \qquad
  Q(x_{j},x_{k})=Q(y_{j},y_{k})=0,
  \qquad
  Q(x_{j},y_{k})=0 \ \ (j\ne k).
\end{equation}
Finally let $\eta\in\bigwedge^{2}V$ be the class with the same coefficients,
\begin{equation}\label{eq:etabasis}
  \eta=\sum_{j=1}^{2n}\epsilon_{j}\,x_{j}\wedge y_{j}.
\end{equation}
\end{setupenv}

Over $\CC$ the operator $S$ has the two eigenvalues $\pm\delta$, and the
eigenvectors are written down at once.

\begin{lemma}[Eigenbasis]\label{lem:eigenbasis}
For $1\le j\le2n$ and $s\in\{+,-\}$ put
\begin{equation}\label{eq:eigenvec}
  u_{j}^{s}\;=\;\tfrac12\Bigl(x_{j}-s\,\tfrac{\delta}{d}\,y_{j}\Bigr)
  \ \in V_{\CC}.
\end{equation}
Then $S u_{j}^{s}=s\,\delta\,u_{j}^{s}$, the vectors $u_{1}^{+},\dots,u_{2n}^{+}$
form a basis of $V_{+}$ and $u_{1}^{-},\dots,u_{2n}^{-}$ a basis of $V_{-}$,
and complex conjugation on $V_{\CC}$ interchanges $u_{j}^{+}$ and $u_{j}^{-}$.
\end{lemma}

\begin{proof}
Using \eqref{eq:Sbasis} and $\delta^{2}=-d$,
\[
  S u_{j}^{s}
  =\tfrac12\bigl(y_{j}+s\,\tfrac{\delta}{d}\,d\,x_{j}\bigr)
  =\tfrac{s\delta}{2}\bigl(x_{j}+\tfrac{1}{s\delta}y_{j}\bigr)
  =\tfrac{s\delta}{2}\bigl(x_{j}-s\tfrac{\delta}{d}y_{j}\bigr)
  =s\,\delta\,u_{j}^{s},
\]
where the third equality is $1/\delta=-\delta/d$. The $2n$ vectors $u_{j}^{+}$
are linearly independent because the change of basis
$(x_{j},y_{j})\mapsto(u_{j}^{+},u_{j}^{-})$ is block diagonal with invertible
$2\times2$ blocks of determinant $\delta/(2d)\ne0$, and the eigenspaces have
dimension $2n$ each by \Cref{not:eigen}. Conjugation fixes $x_{j}$ and $y_{j}$
and sends $\delta$ to $-\delta$, which interchanges the two formulas.
\end{proof}

\begin{definition}[The complex structure]\label{def:modelHS}
Define a Hodge structure of weight one on $V$ by
\begin{equation}\label{eq:V10}
  V^{1,0}
  =\bigl\langle u_{1}^{+},\dots,u_{n}^{+},\,
      u_{n+1}^{-},\dots,u_{2n}^{-}\bigr\rangle_{\CC},
  \qquad
  V^{0,1}=\overline{V^{1,0}} .
\end{equation}
\end{definition}

The choice is the whole point: the first $n$ basis directions contribute to
$V^{1,0}$ through their $+$ eigenvector and the last $n$ through their $-$
eigenvector, so each of $V_{\pm}$ meets $V^{1,0}$ in exactly $n$ dimensions.

\begin{theorem}[The model is of Weil type]\label{thm:model}
With the data of \Cref{setup:model,def:modelHS}, the pair $(V,V^{1,0})$ is a
polarised rational Hodge structure of weight one with polarising form $Q$, so
$V=H^{1}(A,\QQ)$ for an abelian variety $A$ of dimension $2n$; the
multiplication $S$ makes $K$ act on $A$ up to isogeny; $\eta$ is the class of
the polarisation; and $(A,\iota,\eta)$ is of $(K,-1,n)$-Weil type in the sense
of \Cref{def:weiltype}.
\end{theorem}

\begin{proof}
We verify the two Riemann bilinear relations, then the three conditions of
\Cref{def:weiltype}.

\emph{The form vanishes on $V^{1,0}$.} By \eqref{eq:Qbasis} the form $Q$ is
block diagonal for the decomposition of $V$ into the planes
$\langle x_{j},y_{j}\rangle$, so $Q(u_{j}^{s},u_{k}^{t})=0$ whenever $j\ne k$.
Within one block, $Q(u_{j}^{s},u_{j}^{s})=0$ because $Q$ is alternating, and
\begin{equation}\label{eq:Qplusminus}
  Q(u_{j}^{+},u_{j}^{-})
  =\tfrac14\,Q\bigl(x_{j}-\tfrac{\delta}{d}y_{j},\;
                    x_{j}+\tfrac{\delta}{d}y_{j}\bigr)
  =\tfrac14\cdot\tfrac{2\delta}{d}\,Q(x_{j},y_{j})
  =\frac{\delta\,\epsilon_{j}}{2d}.
\end{equation}
The basis \eqref{eq:V10} of $V^{1,0}$ contains, for each $j$, exactly one of
$u_{j}^{+},u_{j}^{-}$. Hence any two of its members either sit in different
blocks or are equal, and in both cases $Q$ vanishes on them. So
$Q(V^{1,0},V^{1,0})=0$.

\emph{Positivity.} Let $0\ne v=\sum_{j}a_{j}v_{j}\in V^{1,0}$, where $v_{j}$ is
$u_{j}^{+}$ for $j\le n$ and $u_{j}^{-}$ for $j>n$. Then
$\bbar{v}=\sum_{j}\bbar{a_{j}}\,\bbar{v_{j}}$ with $\bbar{v_{j}}$ the other
eigenvector in the same block, so again only the diagonal terms survive:
\[
  i\,Q(v,\bbar{v})
  =\sum_{j\le n}i\,|a_{j}|^{2}\,Q(u_{j}^{+},u_{j}^{-})
   +\sum_{j>n}i\,|a_{j}|^{2}\,Q(u_{j}^{-},u_{j}^{+}) .
\]
By \eqref{eq:Qplusminus} and $\delta=i\sqrt{d}$ the first kind of term is
$i\cdot\delta\epsilon_{j}/(2d)=-\epsilon_{j}/(2\sqrt{d})$ and the second is
$+\epsilon_{j}/(2\sqrt{d})$. With $\epsilon_{j}=-1$ for $j\le n$ and $+1$ for
$j>n$ both are $+1/(2\sqrt{d})$, so $i\,Q(v,\bbar{v})=\|a\|^{2}/(2\sqrt{d})>0$.
This is the second Riemann relation, and it is what forces the signature of
$\epsilon$ to be $(n,n)$: reversing any single $\epsilon_{j}$ makes the
corresponding term negative.

\emph{Weil type.} Condition (i) is $\dim A=\tfrac12\dim_{\QQ}V=2n$. For (ii),
$V^{1,0}$ contains $u_{j}^{+}$ for the $n$ indices $j\le n$ and $u_{j}^{-}$ for
the $n$ indices $j>n$, so $\dim V^{1,0}_{+}=\dim V^{1,0}_{-}=n$, which is
\Cref{lem:balanced}, equivalent to the balanced eigenvalue condition on the
tangent space by duality. For (iii), write $\tau=a+b\sqrt{-d}$, so that
$\iota(\tau)$ acts by
\[
  x_{j}\longmapsto a\,x_{j}+b\,y_{j},
  \qquad
  y_{j}\longmapsto -bd\,x_{j}+a\,y_{j},
\]
and $\iota(\bbar{\tau})$ by the same formulas with $b$ replaced by $-b$. Since
$Q$ is block diagonal it suffices to compare on one block, where the three
pairs give
\begin{align*}
  Q\bigl(\iota(\tau)x_{j},y_{j}\bigr)&=a\epsilon_{j}
   =Q\bigl(x_{j},\iota(\bbar{\tau})y_{j}\bigr),\\
  Q\bigl(\iota(\tau)x_{j},x_{j}\bigr)&=-b\epsilon_{j}
   =Q\bigl(x_{j},\iota(\bbar{\tau})x_{j}\bigr),\\
  Q\bigl(\iota(\tau)y_{j},y_{j}\bigr)&=-bd\,\epsilon_{j}
   =Q\bigl(y_{j},\iota(\bbar{\tau})y_{j}\bigr),
\end{align*}
so $\iota(\tau)^{\dagger}=\iota(\bbar{\tau})$. Equivalently, and this is the
form used below, $\iota(\tau)^{*}\eta=\Nm(\tau)\,\eta$, because
$x_{j}\wedge y_{j}=\tfrac{2d}{\delta}\,u_{j}^{+}\wedge u_{j}^{-}$ by
\eqref{eq:eigenvec} and $\iota(\tau)$ multiplies $u_{j}^{+}\wedge u_{j}^{-}$ by
$\tau\bbar{\tau}$.

Finally $\eta$ is of type $(1,1)$, since in each block one of
$u_{j}^{+},u_{j}^{-}$ lies in $V^{1,0}$ and the other in $V^{0,1}$, and
$\eta^{2n}\ne0$, so $\eta$ is a polarisation class.
\end{proof}

\begin{lemma}[The hermitian form]\label{lem:hermform}
On the model of \Cref{setup:model} there is a unique $K$-hermitian form $H$ on
$V$, linear in the first variable, with
\begin{equation}\label{eq:QfromH}
  Q(v,w)=\tfrac12\Tr_{K/\QQ}\bigl(\delta^{-1}H(v,w)\bigr),
  \qquad\text{namely}\qquad
  H(v,w)=Q(\sqrt{-d}\,v,\,w)+\sqrt{-d}\,Q(v,w),
\end{equation}
and in the basis $u_{1},\dots,u_{2n}$ it is
$H=\mathrm{diag}(-\epsilon_{1},\dots,-\epsilon_{2n})$, of signature $(n,n)$.
\end{lemma}

\begin{proof}
$K$-linearity in the first variable is the computation
\[
  H(\sqrt{-d}v,w)=Q(-dv,w)+\sqrt{-d}\,Q(\sqrt{-d}v,w)=\sqrt{-d}\,H(v,w),
\]
and the hermitian symmetry $\overline{H(w,v)}=H(v,w)$ follows from
$Q(v,\sqrt{-d}\,w)=-Q(\sqrt{-d}\,v,w)$, which is condition (iii) of
\Cref{def:weiltype} read backwards. Applying $\delta^{-1}$ and taking the
trace kills the first summand of \eqref{eq:QfromH}, since $Q$ is
$\QQ$-valued and $\Tr_{K/\QQ}(\delta^{-1})=0$, and doubles the second; this
gives $Q$ and also the uniqueness, since $H$ is determined by $Q$ and the
$K$-structure. Finally $H(u_{j},u_{j})=Q(y_{j},x_{j})=-\epsilon_{j}$ and
$H(u_{j},u_{k})=0$ for $j\ne k$, so the signature is the number of indices with
$\epsilon_{j}=-1$ against the number with $\epsilon_{j}=+1$, that is $(n,n)$.
\end{proof}

\begin{remark}\label{rem:twosignatures}
Two signatures have now appeared and they are the same signature seen twice.
The hermitian form $H$ has signature $(n,n)$ by \Cref{lem:hermform}, and the
$K$-action on the tangent space has balanced eigenvalue multiplicities $(n,n)$
by \Cref{def:weiltype}(ii). The Riemann positivity computation in the proof of
\Cref{thm:model} is exactly the passage between them: the sign $\epsilon_{j}$
decides both which eigenvector of the $j$-th block is holomorphic and the sign
of $H$ on that block.
\end{remark}

\begin{remark}\label{rem:lattice}
Taking the lattice $\cO^{2n}$ for an order $\cO\subset K$ and scaling $Q$ to be
integral on it produces an actual abelian variety rather than a Hodge
structure; nothing below depends on the lattice, so we work with the rational
Hodge structure throughout. The $n^{2}$-dimensional family of complex
structures obtained by varying \eqref{eq:V10} over the period domain of the
hermitian form is the Weil family, of dimension $n^{2}$ inside the
$n(2n+1)$-dimensional moduli of principally polarised abelian $2n$-folds; see
\Cref{app:cdk}.
\end{remark}

\subsection{The two generators}

\begin{notation}\label{not:wT}
For $T\subseteq\{1,\dots,2n\}$ write
\[
  w^{(T)}=w^{(T)}_{1}\wedge w^{(T)}_{2}\wedge\cdots\wedge w^{(T)}_{2n},
  \qquad
  w^{(T)}_{j}=
  \begin{cases}
    y_{j}, & j\in T,\\
    x_{j}, & j\notin T .
  \end{cases}
\]
The $2^{2n}$ elements $w^{(T)}$ are part of the monomial basis of
$\bigwedge^{2n}V$, which has dimension $\binom{4n}{2n}$; they are exactly the
monomials containing precisely one of $x_{j},y_{j}$ for every $j$.
\end{notation}

\begin{theorem}[The Weil line in coordinates]\label{thm:explicitgens}
Let $(A,\iota,\eta)$ be the model of \Cref{thm:model} and let
\begin{equation}\label{eq:omegagens}
  \omega_{1}=\sum_{\substack{T\subseteq\{1,\dots,2n\}\\ |T|\ \mathrm{even}}}
    (-1)^{|T|/2}\,d^{\,n-|T|/2}\;w^{(T)},
  \qquad
  \omega_{2}=\sum_{\substack{T\subseteq\{1,\dots,2n\}\\ |T|\ \mathrm{odd}}}
    (-1)^{(|T|-1)/2}\,d^{\,n-(|T|+1)/2}\;w^{(T)} .
\end{equation}
Then
\begin{enumerate}[label=\textup{(\roman*)},nosep]
\item $\omega_{1}$ and $\omega_{2}$ have integer coefficients and each has
  exactly $2^{2n-1}$ nonzero terms;
\item $\omega_{1},\omega_{2}$ is a $\QQ$-basis of the Weil line
  $\HW(A)$ of \Cref{def:weilline};
\item writing $\omega_{\pm}=u_{1}^{\pm}\wedge\cdots\wedge u_{2n}^{\pm}$ for the
  generators of $\bigwedge^{2n}V_{\pm}$,
  \[
    \omega_{+}+\omega_{-}=\frac{\omega_{1}}{2^{2n-1}d^{\,n}},
    \qquad
    \frac{\omega_{+}-\omega_{-}}{\delta}=-\frac{\omega_{2}}{2^{2n-1}d^{\,n}} ;
  \]
\item $\iota(\tau)^{*}\omega_{+}=\tau^{2n}\omega_{+}$ and
  $\iota(\tau)^{*}\omega_{-}=\bbar{\tau}^{\,2n}\omega_{-}$ for all
  $\tau\in\Kx$, so on the plane $\langle\omega_{1},\omega_{2}\rangle$ the
  element $\iota(\tau)$ acts by the $\QQ$-rational matrix of multiplication by
  $\tau^{2n}$ on $K$;
\item $\omega_{1}$ and $\omega_{2}$ are of Hodge type $(n,n)$, hence are Hodge
  classes.
\end{enumerate}
\end{theorem}

\begin{proof}
(iii) and (i). Expanding \eqref{eq:eigenvec} in the product defining
$\omega_{s}$ and collecting the monomials $w^{(T)}$,
\begin{equation}\label{eq:omegaexp}
  \omega_{s}
  =\frac{1}{2^{2n}}\sum_{T}\Bigl(-s\,\frac{\delta}{d}\Bigr)^{|T|}w^{(T)},
\end{equation}
because the factor from index $j$ contributes $x_{j}$ with coefficient $1/2$
and $y_{j}$ with coefficient $-s\delta/(2d)$, and the monomials $w^{(T)}$
appear in the order $j=1,\dots,2n$ with no reordering sign. Therefore
\[
  \omega_{+}+\omega_{-}
  =\frac{1}{2^{2n-1}}\sum_{|T|\ \mathrm{even}}\Bigl(\frac{\delta}{d}\Bigr)^{|T|}w^{(T)},
  \qquad
  \frac{\omega_{+}-\omega_{-}}{\delta}
  =-\frac{1}{2^{2n-1}}\sum_{|T|\ \mathrm{odd}}\frac{\delta^{|T|-1}}{d^{|T|}}\,w^{(T)} ,
\]
the odd and even parts cancelling in the respective sums. Now
$\delta^{2m}=(-d)^{m}$, so for $|T|=2m$ the coefficient in the first sum is
$(-1)^{m}d^{-m}$, and for $|T|=2m+1$ the coefficient in the second is
$-(-1)^{m}d^{-m-1}$. Multiplying by $2^{2n-1}d^{\,n}$ clears denominators,
since $m\le n$ in the first case and $m+1\le n$ in the second, and yields
exactly \eqref{eq:omegagens}. The term counts are
$\sum_{m}\binom{2n}{2m}=\sum_{m}\binom{2n}{2m+1}=2^{2n-1}$.

(ii). The two displayed combinations are fixed by complex conjugation, because
conjugation interchanges $\omega_{+}$ and $\omega_{-}$ by
\Cref{lem:eigenbasis} and sends $\delta$ to $-\delta$; and they visibly have
rational coefficients. They are linearly independent over $\CC$ since
$\omega_{+},\omega_{-}$ are, so they span
$\bigl(\bigwedge^{2n}V_{+}\oplus\bigwedge^{2n}V_{-}\bigr)\cap H^{2n}(A,\QQ)$,
which is $\HW(A)$ by \Cref{def:weilline} and \Cref{prop:weilline}(i).

(iv). The first two identities are immediate, since $\iota(\tau)$ multiplies
each of the $2n$ factors $u_{j}^{+}$ by $\tau$. For the last clause, by (iii)
the complex line $\CC\omega_{+}$ is the $\tau^{2n}$-eigenline inside
$\HW(A)\otimes\CC$, so under the isomorphism $\HW(A)\cong K$ of
\Cref{prop:weilline}(iii) the action of $\iota(\tau)$ is multiplication by
$\tau^{2n}$, whose matrix in the basis $\omega_{1},\omega_{2}$ has rational
entries.

(v). By \eqref{eq:V10} the monomial $\omega_{+}=u_{1}^{+}\wedge\cdots\wedge
u_{2n}^{+}$ has exactly $n$ factors in $V^{1,0}$, namely those with $j\le n$,
and $n$ factors in $V^{0,1}$; so $\omega_{+}$ is of type $(n,n)$, and likewise
$\omega_{-}$. Hence so are the two rational combinations.
\end{proof}

\begin{remark}[The balance is what makes them Hodge classes]\label{rem:balancecount}
Step (v) is the only place the balanced condition enters, and we isolate
it. Suppose we keep everything else and replace \eqref{eq:V10} by a
complex structure for which $\dim V^{1,0}_{+}=p$ and $\dim V^{1,0}_{-}=2n-p$.
Then $\omega_{+}$ has $p$ holomorphic factors and $\omega_{+}$ is of type
$(p,2n-p)$, so $\HW(A)$ is a rational sub-Hodge structure of
$H^{2n}(A,\QQ)$ of type $\{(p,2n-p),(2n-p,p)\}$, and it contains a nonzero
Hodge class if and only if $p=n$. The balanced signature in
\Cref{def:weiltype}(ii) is thus not an extra hypothesis imposed for
convenience; it is exactly the condition under which the Weil line has anything
to do with the Hodge conjecture. This is verified for all $p$ and all
$n\le3$ in \texttt{explicit\_weil.py}, check (W4), and is proved again in
\Cref{app:albert}, \Cref{prop:signatureforced}.
\end{remark}

\subsection{The case \texorpdfstring{$n=2$}{n=2}}

Here $\dim A=4$, $\dim_{\QQ}V=8$, $H^{4}(A,\QQ)$ has dimension
$\binom{8}{4}=70$, and each generator has $2^{3}=8$ terms. Writing the
monomials in the order $j=1,2,3,4$ as in \Cref{not:wT}, \eqref{eq:omegagens}
reads
\begin{align}
  \omega_{1}
  &= d^{2}\,x_{1}x_{2}x_{3}x_{4}
     \;-\;d\!\!\sum_{|T|=2}\!\!w^{(T)}
     \;+\;y_{1}y_{2}y_{3}y_{4},
     \label{eq:n2omega1}\\[2pt]
  \omega_{2}
  &= d\!\!\sum_{|T|=1}\!\!w^{(T)}
     \;-\;\sum_{|T|=3}\!\!w^{(T)} ,
     \label{eq:n2omega2}
\end{align}
with the wedge signs omitted from monomials. For $d=1$, written out in full,
\begin{align*}
  \omega_{1}
  &= x_{1}x_{2}x_{3}x_{4}
   -y_{1}y_{2}x_{3}x_{4}-y_{1}x_{2}y_{3}x_{4}-y_{1}x_{2}x_{3}y_{4}
   -x_{1}y_{2}y_{3}x_{4}-x_{1}y_{2}x_{3}y_{4}-x_{1}x_{2}y_{3}y_{4}\\
  &\qquad{}+y_{1}y_{2}y_{3}y_{4},\\
  \omega_{2}
  &= y_{1}x_{2}x_{3}x_{4}+x_{1}y_{2}x_{3}x_{4}+x_{1}x_{2}y_{3}x_{4}
     +x_{1}x_{2}x_{3}y_{4}\\
  &\qquad{}-y_{1}y_{2}y_{3}x_{4}-y_{1}y_{2}x_{3}y_{4}-y_{1}x_{2}y_{3}y_{4}
     -x_{1}y_{2}y_{3}y_{4} .
\end{align*}
The polarisation power is
$\eta^{2}=2\sum_{j<k}\epsilon_{j}\epsilon_{k}\,x_{j}y_{j}x_{k}y_{k}$,
a sum of six monomials, and these are visibly different monomials from the
sixteen above. That observation is the content of the next subsection.

\subsection{The case \texorpdfstring{$n=3$}{n=3}}

Here $\dim A=6$, $\dim_{\QQ}V=12$, $H^{6}(A,\QQ)$ has dimension
$\binom{12}{6}=924$, and each generator has $2^{5}=32$ terms. The coefficients
depend only on $|T|$, so the two classes are recorded completely by
\Cref{tab:n3coeffs}.

\begin{table}[htbp]
\centering
\renewcommand{\arraystretch}{1.2}
\begin{tabular}{@{}cccc@{\hskip 2.4em}cccc@{}}
\toprule
\multicolumn{4}{c}{$\omega_{1}$} & \multicolumn{4}{c}{$\omega_{2}$}\\
\cmidrule(r{1.6em}){1-4}\cmidrule(l){5-8}
$|T|$ & terms & coefficient & at $d=2$ &
$|T|$ & terms & coefficient & at $d=2$\\
\midrule
$0$ & $1$  & $+d^{3}$ & $+8$  & $1$ & $6$  & $+d^{2}$ & $+4$\\
$2$ & $15$ & $-d^{2}$ & $-4$  & $3$ & $20$ & $-d$     & $-2$\\
$4$ & $15$ & $+d$     & $+2$  & $5$ & $6$  & $+1$     & $+1$\\
$6$ & $1$  & $-1$     & $-1$  &     &      &          &     \\
\midrule
    & $32$ &          &       &     & $32$ &          &     \\
\bottomrule
\end{tabular}
\caption{The Weil classes at $n=3$. Each row gives the number of index sets
$T$ of that size and the common coefficient of the corresponding monomials
$w^{(T)}$ in \eqref{eq:omegagens}. The two classes together involve all $64$
monomials containing exactly one of $x_{j},y_{j}$ for each $j$, out of the
$924$ monomials of $H^{6}$.}
\label{tab:n3coeffs}
\end{table}

For $d=1$ the coefficients are the alternating pattern
$+,-,+,-$ and $+,-,+$, and the two classes are the even and odd halves of
$\prod_{j=1}^{6}(x_{j}-y_{j})$ read as a formal product; for general $d$ the
powers of $d$ record the failure of $\sqrt{-d}$ to be a root of unity, which
is the arithmetic behind \Cref{lem:charsdiffer}.

\begin{figure}[htbp]
\centering
  \includegraphics[width=0.78\textwidth]{fig_weilcube}
\caption{The $2^{2n}$ monomials $w(T)$ of \Cref{not:wT}, at $n=2$, drawn as the
vertices of the $4$-cube in a projection that keeps all sixteen distinct. An
edge joins two monomials whose index sets differ in one element. The colouring
is the bipartition by the parity of $\lvert T\rvert$: the eight even vertices
carry $\omega_{1}$ and the eight odd ones carry $\omega_{2}$, so the two
generators of the Weil line are the two halves of the cube and the coefficient
of a monomial depends only on the distance from the empty set.}
\label{fig:weilcoeff}
\end{figure}



\subsection{An elementary proof that the Weil line misses the divisors}

\Cref{prop:notpoly} was deduced from the character computation. In the model
it has a proof with no representation theory in it at all, and the proof
explains the phenomenon in one line: the Weil classes and the powers of the
polarisation are supported on disjoint sets of monomials.

\begin{proposition}[Disjoint supports]\label{prop:disjoint}
In the model of \Cref{thm:model}, every monomial occurring in $\omega_{1}$ or
in $\omega_{2}$ contains exactly one of $x_{j},y_{j}$ for every
$j\in\{1,\dots,2n\}$, and every monomial occurring in $\eta^{k}$ contains
either both or neither, for every $j$ and every $k\ge0$. Consequently
\[
  \HW(A)\cap\QQ\,\eta^{n}=0 ,
\]
and more generally $\HW(A)$ meets the span of the monomials occurring in
powers of $\eta$ in zero.
\end{proposition}

\begin{proof}
The first assertion is \Cref{not:wT}. For the second, $\eta$ is the sum
\eqref{eq:etabasis} of the $2n$ two-forms $x_{j}\wedge y_{j}$, which pairwise
commute and square to zero, so
\[
  \eta^{k}=k!\sum_{\substack{S\subseteq\{1,\dots,2n\}\\|S|=k}}
    \Bigl(\prod_{j\in S}\epsilon_{j}\Bigr)\,
    \bigwedge_{j\in S}\bigl(x_{j}\wedge y_{j}\bigr),
\]
and the monomial indexed by $S$ contains both $x_{j}$ and $y_{j}$ for $j\in S$
and neither for $j\notin S$. A monomial cannot satisfy both patterns for a
given $j$, so the two families of monomials are disjoint, and $2n\ge1$
guarantees that there is at least one index at which to test. Since
$\omega_{1},\omega_{2}$ span $\HW(A)$ by \Cref{thm:explicitgens}(ii), a nonzero
element of $\HW(A)$ has a nonzero coefficient on some $w^{(T)}$, which does not
occur in $\eta^{n}$.
\end{proof}

\begin{remark}\label{rem:twoproofs}
The two proofs of \Cref{prop:notpoly} measure different things.
\Cref{prop:disjoint} is a statement about one model and a choice of basis, and
it is short because the basis was chosen to make the $K$-action diagonal. The
character proof is basis free and applies to every abelian variety of Weil
type, including those not of the form above, and it identifies what the
obstruction is rather than only that there is one. We record both because the
first is checkable by hand and the second is what \Cref{thm:character}
generalises.
\end{remark}

\subsection{All the Hodge classes}

The coordinates also settle how much of the Hodge conjecture for $A$ is at
stake. Recall that the Hodge classes of an abelian variety are the invariants
of the Hodge group,
$\Hdg^{k}(A)=H^{2k}(A,\QQ)^{\Hg(A)}$, where $\Hg(A)\subseteq\SL(V)$ is the
special Mumford--Tate group.

\begin{lemma}[The Hodge group lands in the special unitary group]%
\label{lem:HgSU}
Let $(A,\iota,\eta)$ be of $(K,-1,n)$-Weil type and let $H$ be the
$K$-hermitian form of \Cref{lem:hermform}. Then
\[
  \Hg(A)\subseteq\SU(V,H) .
\]
\end{lemma}

\begin{proof}
By definition $\Hg(A)$ is a $\QQ$-algebraic subgroup of $\SL(V)$ which fixes
every Hodge class in every tensor construction on $V$. It fixes
$\eta\in\bigwedge^{2}V$, hence preserves $Q$; and it commutes with $\iota(K)$,
because each $\iota(\tau)$ is a Hodge class in $V\otimes V^{\vee}$. A
$\QQ$-linear automorphism of $V$ commuting with the $K$-action and preserving
$Q$ preserves $H$, by the formula \eqref{eq:QfromH} expressing $H$ through $Q$
and the $K$-action. Hence $\Hg(A)\subseteq\mathrm{U}(V,H)$. Now
$\HW(A)$ consists of Hodge classes by \Cref{prop:weilline}(ii), so $\Hg(A)$
fixes $\omega_{+}$; and $g\in\mathrm{U}(V,H)$ acts on
$\bigwedge^{2n}V_{+}$ by $\det_{K}(g)$, so $\det_{K}(g)=1$, which is the
definition of $\SU(V,H)$.
\end{proof}

For a very general member of the Weil family the containment of
\Cref{lem:HgSU} is an equality;%
\footnote{Very general means outside a countable union of proper closed analytic subsets of $\cD_{n,n}$, which is the usual meaning in this subject and is the best that can be asked: the locus where the endomorphism algebra jumps is a countable union of such subsets, by \Cref{thm:cdk}, and it is dense in the classical topology.} this is Weil's computation \cite{Weil77}, in
the form given by Moonen and Zarhin \cite{MZ99}. Granting it, the Hodge classes
can be counted exactly.

\begin{theorem}[The full Hodge ring of a general Weil variety]\label{thm:hdgcount}
Let $(A,\iota,\eta)$ be of $(K,-1,n)$-Weil type with
$\Hg(A)=\SU(V,H)$. Then for every $k$ with $0\le k\le 2n$,
\[
  \dim_{\QQ}\Hdg^{k}(A)=
  \begin{cases}
    3, & k=n,\\[2pt]
    1, & k\ne n,
  \end{cases}
  \qquad\text{and}\qquad
  \Hdg^{n}(A)=\QQ\,\eta^{n}\oplus\HW(A).
\]
\end{theorem}

\begin{proof}
Complexify. The group $\SU(V,H)\otimes\CC$ is $\SL(V_{+})\cong\SL_{2n}(\CC)$,
acting on $V_{+}$ tautologically and on $V_{-}$ through the isomorphism
$V_{-}\cong V_{+}^{\vee}$ induced by $H$. By \Cref{prop:grading},
\[
  H^{2k}(A,\CC)=\bigoplus_{r+s=2k}
    \textstyle\bigwedge^{r}V_{+}\otimes\bigwedge^{s}V_{+}^{\vee},
\]
so it suffices to find the invariants in each summand. An invariant is a
weight-zero vector annihilated by every raising operator. In the weight lattice
of $\SL_{2n}$ the weight of $v_{T}\otimes v_{U}^{\vee}$ is
$\mathbf{1}_{T}-\mathbf{1}_{U}$ modulo $(1,1,\dots,1)$; since the entries of
$\mathbf{1}_{T}-\mathbf{1}_{U}$ lie in $\{-1,0,1\}$, the only multiples of
$(1,\dots,1)$ it can equal are $0$ and $\pm(1,\dots,1)$, giving
\[
  T=U \quad(\text{so } r=s),
  \qquad
  (T,U)=\bigl(\{1,\dots,2n\},\emptyset\bigr),
  \qquad
  (T,U)=\bigl(\emptyset,\{1,\dots,2n\}\bigr).
\]
The last two occur only for $(r,s)=(2n,0)$ and $(0,2n)$, each with a
one-dimensional weight-zero space consisting of the invariant lines
$\bigwedge^{2n}V_{+}$ and $\bigwedge^{2n}V_{-}$, which are $\SL$-invariant
because $\SL$ has trivial determinant. For $r=s$ the weight-zero space has
dimension $\binom{2n}{r}$ and its invariant subspace is the line spanned by the
image of the identity under
$\End(\bigwedge^{r}V_{+})\cong\bigwedge^{r}V_{+}\otimes\bigwedge^{r}V_{+}^{\vee}$,
which is one dimensional because $\bigwedge^{r}V_{+}$ is an irreducible
$\SL_{2n}$-module and Schur's lemma applies. The kernel computation confirming
that the invariant space is exactly this line, for all $r,s$ and all $n\le4$,
is item (VII) of \Cref{sec:verification}.

Collecting: for $k\ne n$ the only contribution is $(r,s)=(k,k)$, giving
dimension one, which is $\QQ\eta^{k}$ since $\eta^{k}$ is a nonzero invariant
lying in that summand by \Cref{lem:NSchar}. For $k=n$ there are three
contributions, $(n,n)$, $(2n,0)$ and $(0,2n)$, of total dimension three; the
first is $\QQ\eta^{n}$ and the last two span $\HW(A)\otimes\CC$ by
\Cref{lem:weilchar}. Descending to $\QQ$ preserves dimensions.
\end{proof}

\begin{corollary}[What the Hodge conjecture asks for here]\label{cor:whatisleft}
Let $(A,\iota,\eta)$ be as in \Cref{thm:hdgcount}. The Hodge conjecture holds
for $A$ in every degree if and only if the single class $\omega_{1}$ of
\eqref{eq:omegagens} is algebraic.
\end{corollary}

\begin{proof}
For $k\ne n$ the Hodge classes are the multiples of $\eta^{k}$, which are
algebraic, being powers of a divisor class. For $k=n$ they are spanned by
$\eta^{n}$, $\omega_{1}$ and $\omega_{2}$ by \Cref{thm:hdgcount}. Finally
$\omega_{2}$ is algebraic as soon as $\omega_{1}$ is: by
\Cref{thm:explicitgens}(iv) the plane $\HW(A)$ is a free $K$-module of rank one
generated by $\omega_{1}$, and the action of $\iota(\tau)$ is induced by an
endomorphism of $A$, hence carries algebraic classes to algebraic classes; some
$\tau$ moves $\omega_{1}$ to a class with a nonzero $\omega_{2}$ component, and
the algebraic classes form a $\QQ$-subspace.
\end{proof}

\begin{remark}\label{rem:onlyoneclass}
\Cref{cor:whatisleft} is the reason the Weil line is the standard test case. On
a general abelian variety of Weil type the Hodge conjecture is not a statement
about a large space of classes with some unknown structure; it is a statement
about one explicitly written class with $2^{2n-1}$ integer coefficients. All
three obstruction theorems of this paper are about the difficulty of producing
that one class, and the construction of \Cref{sec:construction} is an attempt
to produce it.
\end{remark}

\begin{figure}[htbp]
\centering
  \includegraphics[width=0.9\textwidth]{fig_hodgespike}
\caption{The dimensions of $\Hdg^{k}(A)$ for a general abelian variety of
$(K,-1,n)$-Weil type, drawn as a bar field over the two axes $k$ and $n$, at
$n=2,3,4$, as computed in \Cref{thm:hdgcount}. Every degree carries the one
class $\eta^{k}$; only the middle degree $k=n$ carries anything else, and the
excess there is exactly the two dimensions of the Weil line. The profile is
the same for every $n$, which is why the problem does not get structurally
harder as the dimension grows, only larger.}
\label{fig:hodgespike}
\end{figure}
